% Generated draft: requires mathematical and source review.
\documentclass[11pt]{article}
\usepackage[margin=25mm]{geometry}
\usepackage{fontspec}
\setmainfont{Latin Modern Roman}
\usepackage{amsmath,amssymb,mathtools,unicode-math}
\setmathfont{Latin Modern Math}
\usepackage{enumitem,needspace,xurl,hyperref}
\setlength{\parindent}{0pt}
\setlength{\parskip}{5pt}
\setlength{\emergencystretch}{3em}
\newcommand{\N}{\mathbb{N}}
\newcommand{\Z}{\mathbb{Z}}
\newcommand{\Q}{\mathbb{Q}}
\newcommand{\R}{\mathbb{R}}
\newcommand{\C}{\mathbb{C}}
\newcommand{\F}{\mathbb{F}}
\newcommand{\A}{\mathbb{A}}
\newcommand{\PP}{\mathbb{P}}
\newcommand{\E}{\mathbb{E}}
\newcommand{\eps}{\varepsilon}
\newcommand{\dd}{\,\mathrm{d}}
\begin{document}
\section*{UnsolvedMath ID: 30004465}
\textbf{Source problem code:} catalogue entry
\subsection{Definitions and mathematical statement}
A Connes--Chern character relates an appropriate form of \(K\)-theory or \(K\)-homology to cyclic homology or cyclic cohomology. A higher rho invariant is a secondary \(K\)-theoretic invariant associated to suitable geometric operator data. These descriptions identify the terminology only: the supplied record does not specify the character construction, the operator, the algebra, the relevant groups, or the hypotheses required to define the higher rho invariant.

Write, schematically,
\[
\bigl\langle \operatorname{Ch},\rho\bigr\rangle_{\mathrm{an}}
\]
for the analytic expression intended by the source for pairing a Connes--Chern character with a higher rho invariant. This notation is a placeholder for the source-dependent construction; it does not assert that such an expression is defined or convergent for arbitrary character and invariant data.

\textbf{Problem.} Does the analytic expression defining this pairing converge for every admissible instance of the geometric and analytic setting intended in the source report? What additional geometric hypotheses guarantee convergence, and in what further geometric situations does convergence hold?

\textbf{Unresolved formulation.} The record supplies neither the analytic expression nor the class of admissible instances. In particular, it does not identify the cyclic cocycles or character representatives, the algebra and its topology, the operator assumptions, or whether convergence concerns a series, an improper integral, or another limiting construction. Consequently, the universal quantifier above is restricted to the intended but unspecified source setting. A fully quantified mathematical statement requires those definitions and hypotheses; the record does not justify a claim about all conceivable Connes--Chern and higher-rho pairings.

The supplied dated assessment reports convergence under smooth-subalgebra, cocycle-extension, and spectral-gap hypotheses, while reporting that convergence in the general geometric situation remains unknown. These hypothesis names are preserved as supplied qualifications, not as a precise theorem: their exact conditions and their logical relationship are absent from the record.
\subsection{Short English statement}
The question asks whether the analytic formulas pairing Connes--Chern characters with higher rho invariants always converge in the intended source setting, and which further geometric assumptions ensure convergence. The supplied record omits the formulas and admissibility conditions, so a fully precise formulation remains unresolved. Its \texttt{partially\_solved} label and descriptions of known convergence results are supplied curation, not independently verified conclusions.
\subsection{Sources}
\begin{itemize}
\item Supplied local UnsolvedMath registry record, permanent ID \(30004465\), title \emph{Convergence of Connes--Chern and Higher-Rho Pairings}, exact source problem code \texttt{catalogue entry}, in the collection \emph{Oberwolfach Reports Open Problems}. No individual proposer is identified. The record gives proposed year \(2020\).
\item Supplied report citation: \emph{Non-Commutative Geometry and Cyclic Homology (2021)}, Oberwolfach Reports, pp.~919--960, DOI \texttt{10.4171/owr/2020/17}. \url{https://doi.org/10.4171/owr/2020/17}. This is also the supplied literature-source link labelled ``original OWR report.'' The report title year and the record's proposed year are retained without reconciliation.
\item Supplied literature-source link labelled ``geometrically discrete groups theorem'': \url{https://arxiv.org/abs/1811.02086}. The record supplies no authors, paper title, theorem number, or precise hypotheses for this reference.
\item Literature assessment date: \texttt{2026-08-22}. The supplied dated curation labels the problem \texttt{partially\_solved} and reports convergence under smooth-subalgebra, cocycle-extension, and spectral-gap hypotheses, with the general geometric situation unresolved. This assessment and its cited claims have not been independently verified here.
\item The supplied statement status is \texttt{corrected\_verified}. Its verification note reports that the official report was checked and that the phrase ``appearing in the report'' was removed while preserving both convergence questions. This is a supplied verification claim, not a source visit or verification performed for this draft.
\end{itemize}
\end{document}
